% qkit-paper LaTeX preamble for the main paper.
% Loaded via the YAML include-in-header field in index.qmd.

\usepackage{amsthm}
\usepackage{caption}
\usepackage{subcaption}
\usepackage{ragged2e}
\usepackage{setspace}
\usepackage{enumitem}
\usepackage{multirow}
\usepackage{longtable}
\usepackage{tabularx}
\usepackage{float}
\usepackage{sectsty}
\usepackage{orcidlink}
\usepackage{etoolbox}

% Pin footnotes to the bottom of every page, even on short pages
% (default LaTeX would float them up against the body text).
\usepackage[bottom]{footmisc}

% Render title-page \thanks{} (the paper's only footnote unless the
% author adds inline ^[...] notes) as an asterisk instead of the
% article-class default numeric superscript.
% Belt-and-braces: (1) set \thefootnote to symbol form, (2) override
% \@makefnmark so the marker is drawn as a literal * at \normalsize
% (so it does not shrink to invisibility next to the ORCID icon).
\renewcommand{\thefootnote}{\fnsymbol{footnote}}
\makeatletter
\renewcommand\@makefnmark{\hbox{\textsuperscript{\normalsize *}}}
\makeatother

% ---------------------------------------------------------------------
% Caption and float-note styling.
%
% Captions render as a bold "Table 1." label with a period separator,
% then a bold ragged-right title at full text width. \floatnote sets a
% small justified note directly underneath, opening with an italic
% "Note:" -- the layout used by most economics and finance journals.
%
% The [sub] line exists only to stop subcaptions inheriting textfont=bf.
% It deliberately does NOT set labelformat=simple, which would restyle
% subcaptions from "(a) Panel" to "a Panel" and subfigure references
% from 1(a) to 1a. That is subfigure styling, not caption styling.
% ---------------------------------------------------------------------
\captionsetup{font=normal, labelfont=bf, textfont=bf, labelsep=period,
              justification=raggedright, singlelinecheck=false,
              width=1\textwidth, indention=0pt}
\captionsetup[sub]{textfont=normalfont}

% \floatnote{...} -- the note block that sits under a caption.
%
% One definition, used two ways: float-notes.lua appends it inside the
% generated \caption{} for markdown floats, and an author writes it by
% hand immediately after \caption{} in a raw-LaTeX float.
%
% The resets are load-bearing. \parindent must be set AFTER \justifying,
% because ragged2e restores \parindent from \JustifyingParindent and this
% template sets `indent: true` -- with the wrong order the note's first
% line indents and no longer aligns with the caption above it. \leftskip
% and \hangindent clear the caption's hanging indent, which would
% otherwise indent every continuation line.
% The opening \vspace sets the caption-to-note gap, and the arithmetic is
% worth knowing before changing it. The \par that closes the caption uses
% the BODY's baselineskip, which `linestretch: 2` doubles to 28pt, so the
% gap is 28pt + this vspace no matter what font the caption is in.
% Restyling the caption (\captionsetup{font=singlespacing}) does not move
% it -- the interline glue at that \par is already fixed by the outer
% spacing. -6pt therefore yields a 22pt gap: clearly more than the
% note's own 13.6pt leading, so the note reads as a separate block,
% without the hole that \vspace{2pt} (30pt) or \vspace{-4pt} (24pt)
% leaves under the caption. Below about 16pt the note stops reading as a
% block of its own and starts looking like a third caption line.
\newcommand{\floatnote}[1]{%
  \par\vspace{-6pt}%
  \begingroup
    \normalfont\small\singlespacing\justifying
    \leftskip=0pt \rightskip=0pt
    \parindent=0pt \hangindent=0pt \hangafter=0
    \textit{Note:} #1\par
  \endgroup
}

% Bibliography spacing: keep double-spacing off inside the reference
% list even when the paper body uses linestretch: 2, and tighten
% inter-entry space. Pandoc emits references inside CSLReferences (or
% under natbib, the thebibliography list) — this hook applies on each
% environment open, so it stays cheap.
\AtBeginEnvironment{CSLReferences}{%
  \singlespacing%
  \setlength{\parskip}{4pt}%
}
\AtBeginEnvironment{thebibliography}{%
  \singlespacing%
  \setlength{\parskip}{4pt}%
}

% Asymmetric vertical padding around [H]-placed figures. \intextsep
% is used both above and below in-text figures, so we add a small
% positive \vspace before each figure and a small negative \vspace
% after each caption for extra top-heavy breathing room. Values here
% pair well with the section top-space trim below so a figure
% immediately followed by a section heading does not accumulate a
% huge accumulated gap.
\setlength{\intextsep}{20pt plus 2pt minus 2pt}
\setlength{\textfloatsep}{20pt plus 2pt minus 2pt}
\setlength{\floatsep}{20pt plus 2pt minus 2pt}
\BeforeBeginEnvironment{figure}{\vspace{10pt}}
\AfterEndEnvironment{figure}{\vspace{-10pt}}

% Keep tables tighter than the double-spaced body. Pandoc renders
% markdown pipe tables as `longtable`, NOT `tabular`, so hooking only
% `tabular` leaves every markdown table inheriting linestretch and
% rendering double-spaced. Both environments are hooked so pipe tables
% and knitr/raw tabular tables look the same.
%
% For a document with long tables, \singlespacing here is often better:
% at \onehalfspacing a 15-row table can overflow a page, and because
% longtable breaks across pages rather than floating, it will split and
% orphan rows rather than move to the next page.
\AtBeginEnvironment{tabular}{\onehalfspacing}
\AtBeginEnvironment{longtable}{\onehalfspacing}

% (The "Preliminary draft" subtitle and the author block with ORCID +
% \thanks{} footnotes are emitted by the custom title.tex partial
% activated via "template-partials: title.tex" in the YAML header.)

% (The author-separator logic lives in title.tex and is driven by
% the is-first / is-last / is-only-coauthor flags injected onto the
% YAML author list by mark-last-author.lua. No preamble macros are
% needed.)

% Italic subsection / subsubsection headings, with both sectsty and
% an explicit \@startsection override so the article-class default
% \bfseries cannot leak through. Trim the "before" space on section /
% subsection headings so a figure immediately followed by a heading
% does not accumulate a huge gap (paired with the \intextsep and
% \BeforeBeginEnvironment{figure} hooks above).
\subsectionfont{\large\itshape}
\subsubsectionfont{\large\itshape}
\makeatletter
\renewcommand{\section}{\@startsection{section}{1}{\z@}%
  {-1.5ex \@plus -0.5ex \@minus -.2ex}{1.0ex \@plus .2ex}%
  {\normalfont\Large\bfseries}}
\renewcommand{\subsection}{\@startsection{subsection}{2}{\z@}%
  {-1.5ex \@plus -0.5ex \@minus -.2ex}{0.8ex \@plus .2ex}%
  {\normalfont\large\itshape}}
\renewcommand{\subsubsection}{\@startsection{subsubsection}{3}{\z@}%
  {-1.5ex \@plus -0.5ex \@minus -.2ex}{0.8ex \@plus .2ex}%
  {\normalfont\large\itshape}}
\makeatother

% Period after the section number.
\makeatletter
\def\@seccntformat#1{\@ifundefined{#1@cntformat}%
    {\csname the#1\endcsname\quad}%
    {\csname #1@cntformat\endcsname}%
}
\def\section@cntformat{\thesection.\quad}
\def\subsection@cntformat{\thesubsection \quad}
\makeatother

\tolerance=1
\emergencystretch=\maxdimen
\hyphenpenalty=10000
\hbadness=10000

\renewenvironment{quote}
  {\normalsize\list{}{\rightmargin=1.5cm \leftmargin=1.5cm}%
   \item\relax}
  {\endlist}

% ---------------------------------------------------------------------
% Symmetric display-equation spacing.
%
% Two independent problems, and BOTH have to be fixed or the gap above a
% display will not match the gap below it.
%
% (1) The four display-skip lengths. The short-skip variants kick in when
% the line immediately preceding the display is short, and by
% article-class default \abovedisplayshortskip is 0pt while
% \belowdisplayshortskip is 7pt; setspace then rescales all four by the
% linestretch. \qkitdisplayskips pins them to one common value, applied
% at load time, again at \AtBeginDocument (after setspace has had its
% say), and once more at every display open via \everydisplay and the
% per-environment hooks below, so amsmath cannot reset them locally.
%
% (2) The empty paragraph above the display -- the bigger of the two, and
% the reason (1) alone never made the spacing look symmetric. Pandoc
% always emits a display as its own block, separated from the prose by a
% blank line, so LaTeX meets \begin{equation} in VERTICAL mode. TeX opens
% a paragraph, finds it empty, and contributes a full \baselineskip of
% interline glue -- 28.8pt at `linestretch: 2` -- on top of
% \abovedisplayskip. Nothing below the display matches it, so the gap
% above was ~29pt larger no matter how the four skips were set.
%
% \qkitdisplayfix cancels exactly that one \baselineskip. It fires only
% in vertical mode, so a display that genuinely continues a paragraph
% (no blank line before it, as in hand-written LaTeX) is left alone, and
% only when \prevdepth shows interline glue will actually be inserted --
% at the top of a page there is no empty-paragraph glue to cancel.
% The cancellation is exact: the previous line's depth appears on both
% sides of the arithmetic and drops out, so the correction does not
% depend on whether that line happened to end in a descender.
% ---------------------------------------------------------------------
\newlength{\qkitdisplayskip}
\setlength{\qkitdisplayskip}{12pt}
\newcommand{\qkitdisplayskips}{%
  \setlength{\abovedisplayskip}{\qkitdisplayskip}%
  \setlength{\belowdisplayskip}{\qkitdisplayskip}%
  \setlength{\abovedisplayshortskip}{\qkitdisplayskip}%
  \setlength{\belowdisplayshortskip}{\qkitdisplayskip}%
}
\makeatletter
\newcommand{\qkitdisplayfix}{%
  \qkitdisplayskips
  \ifvmode\ifdim\prevdepth>-\@m pt \vskip-\baselineskip\fi\fi
}
\makeatother
\qkitdisplayskips
\AtBeginDocument{\qkitdisplayskips}
\everydisplay=\expandafter{\the\everydisplay \qkitdisplayskips}
% amsmath defines \[ ... \] as the equation* environment, so the
% equation* hook also covers unlabelled `$$ ... $$` displays.
\BeforeBeginEnvironment{equation}{\qkitdisplayfix}
\BeforeBeginEnvironment{equation*}{\qkitdisplayfix}
\BeforeBeginEnvironment{align}{\qkitdisplayfix}
\BeforeBeginEnvironment{align*}{\qkitdisplayfix}
\BeforeBeginEnvironment{gather}{\qkitdisplayfix}
\BeforeBeginEnvironment{gather*}{\qkitdisplayfix}
\BeforeBeginEnvironment{multline}{\qkitdisplayfix}
\BeforeBeginEnvironment{multline*}{\qkitdisplayfix}

\theoremstyle{plain}
\newtheorem{theorem}{Theorem}[section]
\newtheorem{proposition}[theorem]{Proposition}
\newtheorem{lemma}[theorem]{Lemma}
\theoremstyle{definition}
\newtheorem{definition}[theorem]{Definition}
\newtheorem{assumption}[theorem]{Assumption}

% ---------------------------------------------------------------------
% Citation styling (natbib).
%
% natbib emits an author-year citation as \NAT@nmfmt{<authors>}, then
% \hyper@natlinkbreak, then \NAT@date -- all inside one \NAT@hyper@
% group. hyperref turns that into two adjacent link annotations, so the
% author names are clickable too and both halves pick up citecolor.
% Quarto's default citeproc instead leaves the names outside the link
% and tints only the year.
%
% \setcitestyle{aysep={}} drops natbib's comma before the year, giving
% "(Campbell and Shiller 1988)" in the Chicago author-date form rather
% than "(Campbell and Shiller, 1988)".
% ---------------------------------------------------------------------
\setcitestyle{aysep={}}
